\section{Candidate statement and scope}
\label{sec:p147-candidate-scope}

This is an isolated research candidate, not an accepted manuscript theorem.
Independent mathematical review and the unchanged production permission and
review gates are required before any integration.

We use squared Hermitian norms and
\(\mathrm d\mathrm d^{\mathrm c}
 =\frac{\mathrm i}{2\pi}\partial\bar\partial\).
Let \(X\) be a connected smooth complex projective variety, let \(L\) be a
big line bundle, let \(h\) be a smooth Hermitian metric, and put
\(\theta=c_1(L,h)\).  Set
\[
 V_\theta=\sup\{v\in\operatorname{PSH}(X,\theta):v\leq0\}.
\]
Write \(v\preceq u\) if \(v\leq u+C\) for some real constant \(C\), and use
the normalized singularity envelope
\[
 P_\theta[u]
 =\left(\sup\{v\in\operatorname{PSH}(X,\theta):
                  v\leq0,\ v\preceq u\}\right)^*.
\]
Here \({}^*\) denotes upper semicontinuous regularization.  The obstacle
\(0\) gives the same envelope as \(V_\theta\), since every
\(\theta\)-psh function below \(0\) is below \(V_\theta\).
The multiplier-ideal envelope \(P_\theta[u]_{\mathcal I}\) is the
normalized envelope with the conditions
\(\mathcal I(tv)\subseteq\mathcal I(tu)\) for every \(t\geq0\).

Let \(\mathcal F\) be a decreasing, left-continuous, multiplicative,
linearly bounded filtration of \(R(X,L)\), precisely as in
\cite[Definition~4.1]{DXZ23}.  Put
\[
 c=\tau_L(\mathcal F),\qquad
 W_m^\tau=\mathcal F^{m\tau}H^0(X,L^{\otimes m}),
\]
and, for \(\tau<c\), use the exact normalization of
\cite[Definition~4.3]{DXZ23}:
\begin{equation}\label{eq:p147-raw-definition}
 u_{\tau,m}
 =\sup_{\substack{s\in W_m^\tau\\
             \sup_X|s|_{h^m}^2\leq1}}
       \frac1m\log|s|_{h^m}^2,\qquad
 u_\tau=\left(\sup_{m\geq1}u_{\tau,m}\right)^*.
\end{equation}
An empty nonzero section space gives \(-\infty\).
For a nonzero finite-dimensional space the first supremum already gives
a \(\theta\)-psh function with an analytic singularity type.
The factor is \(1/m\), not \(1/(2m)\).

\begin{theorem}\label{thm:p147-filtration-maximality-candidate}
With these hypotheses, for every \(\tau<c\),
\begin{equation}\label{eq:p147-target}
 u_\tau=P_\theta[u_\tau]=P_\theta[u_\tau]_{\mathcal I}.
\end{equation}
In particular the raw filtration test curve is maximal on its entire
strict interior, which is the positive-mass range in the source.
\end{theorem}

The proof below addresses the ordinary equality \(u_\tau=P_\theta[u_\tau]\).
The second equality is the already published
\cite[Lemma~4.5]{DXZ23}, not a separate conjecture.
No finite-generation assumption on \(\mathcal F\) is imposed.

\begin{corollary}\label{cor:p147-terminal-slice-candidate}
If the particular DXZ terminal slice
\(u_c=\lim_{\tau\uparrow c}u_\tau\) is not identically \(-\infty\), then
\[
                 u_c=P_\theta[u_c]=P_\theta[u_c]_{\mathcal I}.
\]
This holds whether its top nonpluripolar mass is positive or zero.
\end{corollary}

The same proof does not assert recovery of an arbitrary zero-mass
potential from its ideals, a comparison with the MP cone, or a theory for
arbitrary transcendental pseudoeffective classes.
A separate elementary endpoint observation appears after the proof; it
uses only the particular left-limit convention of DXZ and is not a
general zero-mass reconstruction theorem.
